% Generated draft: requires mathematical and source review.
\documentclass[11pt]{article}
\usepackage[margin=25mm]{geometry}
\usepackage{fontspec}
\setmainfont{Latin Modern Roman}
\usepackage{amsmath,amssymb,mathtools,unicode-math}
\setmathfont{Latin Modern Math}
\usepackage{enumitem,needspace,xurl,hyperref}
\setlength{\parindent}{0pt}
\setlength{\parskip}{5pt}
\setlength{\emergencystretch}{3em}
\newcommand{\N}{\mathbb{N}}
\newcommand{\Z}{\mathbb{Z}}
\newcommand{\Q}{\mathbb{Q}}
\newcommand{\R}{\mathbb{R}}
\newcommand{\C}{\mathbb{C}}
\newcommand{\F}{\mathbb{F}}
\newcommand{\A}{\mathbb{A}}
\newcommand{\PP}{\mathbb{P}}
\newcommand{\E}{\mathbb{E}}
\newcommand{\eps}{\varepsilon}
\newcommand{\dd}{\,\mathrm{d}}
\begin{document}
\section*{UnsolvedMath ID: 2567}
\textbf{Source problem code:} catalogue entry
\subsection{Definitions and mathematical statement}
For subsets $X,Y$ of a group $G$, let $XY=\{xy:x\in X,y\in Y\}$. The product is direct if the multiplication map $X\times Y\to G$ is injective, equivalently every element of $XY$ has a unique representation $xy$ with $x\in X$ and $y\in Y$. Directness does not require $XY=G$.

For a fixed $A\subseteq G$, let $\mathcal D_A=\{B\subseteq G:AB\text{ is direct}\}$. A member $B$ is maximal if there is no strictly larger subset $B'\supsetneq B$ in $\mathcal D_A$. This means maximal by inclusion, not maximum cardinality. The question is whether
\[
\exists\text{ infinite group }G\ \forall A\subseteq G\ \forall B,C\in\mathcal D_A,
\]
\[
\bigl(B,C\text{ are inclusion-maximal in }\mathcal D_A\bigr)
\ \Longrightarrow\ |B|=|C|.
\]
All subsets, not only subgroups or generating sets, are quantified.

For nonempty $A$ and a maximal $B$, also nonempty, the supplied discussion records the equivalent conditions
\[
A^{-1}A\cap BB^{-1}=\{1\},\qquad A^{-1}AB=G,
\]
where $1$ is the group identity and inverse/product notation is setwise. The first expresses directness and the second inclusion-maximality. The discussion attributed to M. H. Hooshmand uses cardinal arithmetic to obtain $|G|=|A||B|$ and reduces a potential verification to balanced subsets with $|A|=|G\setminus A|$: if $|A|<|G\setminus A|$, every maximal $B$ has size $|G|$; if $|A|>|G\setminus A|$, one has $A^{-1}A=G$ and every maximal $B$ has size one. The empty case is separate: if $A=\varnothing$, every $AB$ is vacuously direct and the sole inclusion-maximal $B$ is $G$, so it causes no obstruction.

The supplied solution, attributed to Kabenyuk's 2026 preprint, reports the exact negative statement
\[
\forall\text{ infinite groups }G\ \exists\varnothing\ne A\subseteq G\ \exists B,C\subseteq G,
\]
\[
B,C\text{ are maximal in }\mathcal D_A\quad\text{and}\quad |B|\ne|C|.
\]
It describes every infinite group as unstable in this direct-complement sense. This gives a nonempty witness and refutes the universal-subset property despite the permitted empty subset.

The registry and literature triage dated 2026-08-17 classify the existence question as solved negatively. The source retains peer review or publication of the recent exact-match preprint as a follow-up and qualifies the conclusion by the correctness of its theorem. These are supplied result and status reports, not an independent proof review in this draft.
\subsection{Short English statement}
Can an infinite group have the same cardinality for every inclusion-maximal direct right complement of each fixed subset? The supplied 2026 preprint says no: every infinite group has a nonempty subset with two such maximal complements of different cardinalities.
\subsection{Sources}
\begin{itemize}
\item Kourovka Notebook, New Problems, 21st issue (2026), original Problem 21.58; supplied reduction discussion attributed to M. H. Hooshmand.
\item Mikhail Kabenyuk, ``Factors in infinite groups,'' arXiv:2602.22876 (2026): \url{https://arxiv.org/abs/2602.22876}. The supplied discussion credits M. I. Kabenyuk with the negative solution.
\item Local registry data/unsolvedmath/problems.json, permanent catalogue entry, source code catalogue entry, ``Kourovka Notebook Problem 21.58''; statement, complete supplied discussion and preprint attribution, and literature triage dated 2026-08-17. Dataset attribution: Ulam AI and the UnsolvedMath Contributors.
\end{itemize}
\end{document}
